\begin{table}[H]
\centering
\caption{Tarea de ingeniería 1: Implementación del menú y pantalla de instrucciones (elaboración propia).\label{tab:ti1}}
\begin{tabular}{|p{0.3\textwidth}|p{0.65\textwidth}|}
\hline
\textbf{Campo} & \textbf{Descripción} \\
\hline
Tarea \# & 1 \\
\hline
HU relacionada & HU \#1, HU \#2 \\
\hline
Nombre & Implementación del menú principal y pantalla de instrucciones \\
\hline
Tipo & Desarrollo \\
\hline
Puntos estimados & 2.0 días \\
\hline
Iteración & 1 \\
\hline
Programador & Ruslan Alexei Martinez Campos \\
\hline
Descripción & Se diseñan las interfaces gráficas del menú principal y la pantalla de instrucciones en Godot, incluyendo navegación y elementos educativos. \\
\hline
\end{tabular}
\end{table}

\begin{table}[H]
\centering
\caption{Tarea de ingeniería 2: Visualización e inicio de partida (elaboración propia).\label{tab:ti2}}
\begin{tabular}{|p{0.3\textwidth}|p{0.65\textwidth}|}
\hline
\textbf{Campo} & \textbf{Descripción} \\
\hline
Tarea \# & 2 \\
\hline
HU relacionada & HU \#3, HU \#4 \\
\hline
Nombre & Visualización e inicio de partida \\
\hline
Tipo & Desarrollo \\
\hline
Puntos estimados & 7.0 días \\
\hline
Iteración & 1 \\
\hline
Programador & Ruslan Alexei Martinez Campos \\
\hline
Descripción & Se modela el tablero hexagonal de 63 casillas y se implementa la interfaz de configuración de partida para instanciar a los jugadores. \\
\hline
\end{tabular}
\end{table}

\begin{table}[H]
\centering
\caption{Tarea de ingeniería 3: Lógica del dado y movimiento de fichas (elaboración propia).\label{tab:ti3}}
\begin{tabular}{|p{0.3\textwidth}|p{0.65\textwidth}|}
\hline
\textbf{Campo} & \textbf{Descripción} \\
\hline
Tarea \# & 3 \\
\hline
HU relacionada & HU \#5, HU \#6 \\
\hline
Nombre & Implementación de la lógica del dado y movimiento animado \\
\hline
Tipo & Desarrollo \\
\hline
Puntos estimados & 5.5 días \\
\hline
Iteración & 2 \\
\hline
Programador & Ruslan Alexei Martinez Campos \\
\hline
Descripción & Se programa la lógica pseudoaleatoria del dado y se implementa el controlador para animar el recorrido de las fichas casilla a casilla. \\
\hline
\end{tabular}
\end{table}

\begin{table}[H]
\centering
\caption{Tarea de ingeniería 4: Lógica de efectos de casillas especiales (elaboración propia).\label{tab:ti4}}
\begin{tabular}{|p{0.3\textwidth}|p{0.65\textwidth}|}
\hline
\textbf{Campo} & \textbf{Descripción} \\
\hline
Tarea \# & 4 \\
\hline
HU relacionada & HU \#7, HU \#8 \\
\hline
Nombre & Lógica de efectos de casillas especiales \\
\hline
Tipo & Desarrollo \\
\hline
Puntos estimados & 3.0 días \\
\hline
Iteración & 2 \\
\hline
Programador & Ruslan Alexei Martinez Campos \\
\hline
Descripción & Se desarrolla la lógica de aterrizaje para aplicar desplazamientos dinámicos al detectar colisiones con casillas de tipo escalera o serpiente. \\
\hline
\end{tabular}
\end{table}

\begin{table}[H]
\centering
\caption{Tarea de ingeniería 5: Lógica de preguntas y evaluación educativa (elaboración propia).\label{tab:ti5}}
\begin{tabular}{|p{0.3\textwidth}|p{0.65\textwidth}|}
\hline
\textbf{Campo} & \textbf{Descripción} \\
\hline
Tarea \# & 5 \\
\hline
HU relacionada & HU \#9, HU \#10 \\
\hline
Nombre & Lógica de preguntas y evaluación educativa \\
\hline
Tipo & Desarrollo \\
\hline
Puntos estimados & 8.0 días \\
\hline
Iteración & 3 \\
\hline
Programador & Ruslan Alexei Martinez Campos \\
\hline
Descripción & Se integra el motor de lectura JSON para cargar preguntas ODS y se programa el panel modal interactivo para evaluar y retroalimentar al jugador. \\
\hline
\end{tabular}
\end{table}

\begin{table}[H]
\centering
\caption{Tarea de ingeniería 6: Gestión de turnos del juego (elaboración propia).\label{tab:ti6}}
\begin{tabular}{|p{0.3\textwidth}|p{0.65\textwidth}|}
\hline
\textbf{Campo} & \textbf{Descripción} \\
\hline
Tarea \# & 6 \\
\hline
HU relacionada & HU \#11 \\
\hline
Nombre & Gestión de turnos del juego \\
\hline
Tipo & Desarrollo \\
\hline
Puntos estimados & 2.0 días \\
\hline
Iteración & 3 \\
\hline
Programador & Ruslan Alexei Martinez Campos \\
\hline
Descripción & Se programa el bucle principal de juego en GameManager para orquestar la concurrencia de turnos y bloquear interacciones asíncronas. \\
\hline
\end{tabular}
\end{table}

\begin{table}[H]
\centering
\caption{Tarea de ingeniería 7: Implementación de condición matemática de victoria y podio (elaboración propia).\label{tab:ti7}}
\begin{tabular}{|p{0.3\textwidth}|p{0.65\textwidth}|}
\hline
\textbf{Campo} & \textbf{Descripción} \\
\hline
Tarea \# & 7 \\
\hline
HU relacionada & HU \#12, HU \#13 \\
\hline
Nombre & Implementación de condición matemática de victoria y podio \\
\hline
Tipo & Desarrollo \\
\hline
Puntos estimados & 5.5 días \\
\hline
Iteración & 4 \\
\hline
Programador & Ruslan Alexei Martinez Campos \\
\hline
Descripción & Se programa la condición de victoria en la casilla 63 y se implementa el gestor de puntuaciones para generar la pantalla de estadísticas finales. \\
\hline
\end{tabular}
\end{table}

\begin{table}[H]
\centering
\caption{Tarea de ingeniería 8: Configuración de controles de partida y audio (elaboración propia).\label{tab:ti8}}
\begin{tabular}{|p{0.3\textwidth}|p{0.65\textwidth}|}
\hline
\textbf{Campo} & \textbf{Descripción} \\
\hline
Tarea \# & 8 \\
\hline
HU relacionada & HU \#14, HU \#15, HU \#16 \\
\hline
Nombre & Configuración de controles de partida y audio \\
\hline
Tipo & Desarrollo \\
\hline
Puntos estimados & 4.5 días \\
\hline
Iteración & 4 \\
\hline
Programador & Ruslan Alexei Martinez Campos \\
\hline
Descripción & Se ensambla el menú de pausa in-game conectando los controladores globales de volumen de audio y la lógica de reinicio de estado de la partida. \\
\hline
\end{tabular}
\end{table}

\begin{table}[H]
\centering
\caption{Tarea de ingeniería 9: Implementación de la cámara de seguimiento dinámico (elaboración propia).\label{tab:ti9}}
\begin{tabular}{|p{0.3\textwidth}|p{0.65\textwidth}|}
\hline
\textbf{Campo} & \textbf{Descripción} \\
\hline
Tarea \# & 9 \\
\hline
HU relacionada & HU \#17 \\
\hline
Nombre & Implementación de la cámara de seguimiento dinámico \\
\hline
Tipo & Desarrollo \\
\hline
Puntos estimados & 2.0 días \\
\hline
Iteración & 5 \\
\hline
Programador & Ruslan Alexei Martinez Campos \\
\hline
Descripción & Se instancia una Camera2D controlada mediante interpolación para seguir dinámicamente la trayectoria de la ficha del jugador activo. \\
\hline
\end{tabular}
\end{table}

\begin{table}[H]
\centering
\caption{Tarea de ingeniería 10: Implementación del modo de práctica en solitario (elaboración propia).\label{tab:ti10}}
\begin{tabular}{|p{0.3\textwidth}|p{0.65\textwidth}|}
\hline
\textbf{Campo} & \textbf{Descripción} \\
\hline
Tarea \# & 10 \\
\hline
HU relacionada & HU \#18 \\
\hline
Nombre & Implementación del modo de práctica en solitario \\
\hline
Tipo & Desarrollo \\
\hline
Puntos estimados & 5.0 días \\
\hline
Iteración & 5 \\
\hline
Programador & Ruslan Alexei Martinez Campos \\
\hline
Descripción & Se adapta el controlador de turnos para omitir restricciones de cambio de jugador cuando se inicializa la partida en modo de práctica solitario. \\
\hline
\end{tabular}
\end{table}
